\documentclass[../../main.tex]{subfiles}
\begin{document}

\begin{center}
  {\fontsize{48pt}{58pt}\selectfont \textbf{1996}\normalsize\textbf{年度}}
\end{center}

\begin{table}[htb]
  \centering
  \caption{出題テーマと難易度}
  \begin{tabular}{|c|p{9cm}|c|}
    \hline
    \textbf{問題番号} & \textbf{テーマ} & \textbf{難易度} \\
    \hline
    第1問 & 一次変換による円の像と接線 & \\
    \hline
    第2問 & 連立不等式から定まる二次方程式の解 & \\
    \hline
    第3問 & 球面上の点から見える立方体の頂点 & \\
    \hline
    第4問 & サイコロの目の記録と期待値 & \\
    \hline
    第5問 & 水面の上昇と回転体の半径 & \\
    \hline
    第6問 & 楕円が円板に含まれる条件と面積 & \\
    \hline
  \end{tabular}
\end{table}

\begin{table}[htb]
  \centering
  \caption{解答}
  \renewcommand{\arraystretch}{1.4}
  \begin{tabular}{|c|c|p{9cm}|}
    \hline
    \textbf{問題番号} & \textbf{小問} & \textbf{答え} \\
    \hline
    第1問 & -- & \\
    \hline
    第2問 & -- & 二次方程式は$-1<x<1$に異なる2実根をもつ
    （証明問題） \\
    \hline
    第3問 & (1) & $0<r\le\dfrac{l}{2}$ \\
    \cline{2-3}
    & (2) & $0<r\le\dfrac{\sqrt{3}}{6}l$ \\
    \hline
    第4問 & (1) & $\displaystyle\lim_{n\to\infty}E(n)=1$ \\
    \cline{2-3}
    & (2) & $\displaystyle\lim_{n\to\infty}N(n)=\dfrac{1}{4}$ \\
    \hline
    第5問 & -- & $r(z)=R\sqrt{\dfrac{z}{z+1}}$ \\
    \hline
    第6問 & (1) & $0<\alpha\le\dfrac{1}{2}$なら
    $0<\beta\le\dfrac{1}{4}$，
    $\dfrac{1}{2}\le\alpha<1$なら
    $0<\beta\le\alpha-\alpha^2$（図示は解答参照） \\
    \cline{2-3}
    & (2) & 面積の最大値は$\dfrac{2\sqrt{3}}{9}\pi$ \\
    \hline
  \end{tabular}
\end{table}

\paragraph{試験形態}
6問．試験時間は要確認．

\end{document}
